\section{Discussion}
\label{sec:discussion}

\paragraph{Deployment guidance.}
Within the shifts, source city, and masking model evaluated here, the results
suggest several practical choices.
Under silent-failure zero-mask dropout, DCRNN, AGCRN, and STAEformer degrade
much less than GWNet; the GWNet gap is tied to dense softmax routing and the
mask-aware variant recovers 75\% without retraining. For zero-shot cross-city
transfer on METR-LA/PEMS-BAY-style highway networks, MegaCRN and DCRNN
perform well. Weekday/weekend comparisons require an explicit metric: MegaCRN and
STAEformer have positive normalized gaps but negative masked-mph gaps.
Their normalized values alone should not drive temporal deployment guidance.

\paragraph{Implications for robustness evaluation.}
In our bank, the clean-test ranking changes under sensor dropout and only
partly predicts cross-city transfer. Richer relation learning is not uniformly
more robust; weaknesses appear on different axes. Benchmarks should therefore
report per-shift results alongside in-distribution accuracy. The GWNet controls
make simple optimizer failure an unlikely explanation for its fragility
(Table~\ref{tab:consolidated}).

\paragraph{Recommendation artifact.}
The released $2{\times}2{\times}2$ decision table maps sensor reliability,
city transfer, and weekend exposure to observed fragility and mitigations
(Appendix~\ref{sec:appendix_decision_tree}). It covers the eight-model bank
and demotes MegaCRN's smoothing-related negative sensor delta. Its temporal
axis uses normalized gaps. AGCRN, recommended for two weekend-exposed
profiles, has the only positive masked-mph weekend gap among the extended
models (Table~\ref{tab:temporal_absolute}); these cells should not be used to
choose models for physical-error performance.

\paragraph{Limitations.}
All datasets are highway loop-detector data. METR-LA is the main training
source, with PEMS-BAY and PEMS04 used for targeted replications; urban
arterials and probe data remain untested. The model set is not exhaustive
(e.g., ASTGCN~\citep{guo2019astgcn}, DSTAGNN~\citep{lan2022dstagnn}), and the
main results use three training seeds. Zero-shot transfer is a diagnostic
floor, not a deployment recommendation; the few-shot probe covers four
backbones and limited adaptation (Appendix~\ref{sec:appendix_fewshot}). The
weekday/weekend comparison is a subgroup analysis within one test split,
not an OOD test. Walk-forward evaluation remains future work, and the gap's
sign depends on normalization and target masking. The six-architecture
concentration correlation weakens on PEMS-BAY and PEMS04; interventions
support the link only for GWNet. MegaCRN's negative sensor delta reflects
adjacency-mediated smoothing, and TOD/DOW covariates alone do not explain
its normalized temporal gap.

\paragraph{Scope of the zero-mask diagnostic.}
Zero is a physically possible speed, so a zero reading alone cannot identify
failure. Corruptions are assigned by a failure mask, not by observed values;
valid zeros at kept sensors are not classified as failures. Both remedies
require identifying failed sensors: LOCF uses the true failure mask
(Appendix~\ref{sec:appendix_imputation}), and detector errors increase MAE for
mask-aware routing (Appendix~\ref{sec:appendix_mask_fpfn}). We did not test LOCF
under imperfect detection. The gap applies where failures go undetected,
telemetry is delayed, or carry-forward is unreliable.
